\documentclass{situationist-preprint}

\title{The Status of Unreviewed Results}
\msid{SCI-01-c}
\msfamily{SCI-01}
\msdiscipline{Philosophy of Science}
\mstype{Consequence}
\mskeywords{preprints; peer review; demarcation; error-correction; scientific record}

\begin{document}
\maketitle

\begin{abstract}
Is a result that has not passed peer review scientific? The question has become practical with the growth of preprint servers, on which results circulate publicly before review and are sometimes cited and acted on without it. We derive an answer from the conclusion of SCI-01, that science is distinguished by institutionalized error-correction. A result is scientific in the derivative sense of SCI-01-a when it is exposed to the standing checks of a scientific enterprise. A private result is not yet so exposed; a public preprint is, although it is not yet part of the record. Peer review is one check among several, and it is neither necessary nor sufficient for a result's scientific status.
\end{abstract}

\section{The question}\label{s:question}

Three kinds of result lack peer review. Some are private: Cavendish's electrical experiments, unpublished in his lifetime, are the example of SCI-01-a. Some are public but unreviewed, such as preprints posted to servers like arXiv, which has carried physics and mathematics in this form since 1991 \citep{Ginsparg2011}. And some are public, unreviewed, and never submitted for review at all. A common view holds that review confers scientific status, so that all three kinds are pre-scientific. The criterion of SCI-01 implies a different classification.

Recall the framework of SCI-01-a. By Principle~P1 the unit of demarcation is the enterprise, and a claim is scientific in a derivative sense by belonging to the record of a scientific enterprise. By Definition~3.1 the record is the set of claims the enterprise treats as established, and by Theorem~4.4 an enterprise is a science if and only if its checks are standing. We add one notion.

\begin{definition}[Candidate]\label{d:candidate}
A claim is a \emph{candidate} of an enterprise at a time if it has been made available to the members of the enterprise in a form that allows them to check it, but has not been admitted to the record.
\end{definition}

\section{Consequences}\label{s:consequences}

\begin{proposition}[Exposure]\label{p:exposure}
A candidate of a scientific enterprise is scientific in the derivative sense: it is subject to the enterprise's standing checks. A private result is not, whatever its quality.
\end{proposition}

\begin{proof}
What makes an enterprise scientific, by Theorem~4.4 of SCI-01-a, is that its checks recur without a last occasion. A check, by Definition~3.2 of SCI-01-a, is a procedure by which members other than the producer test a claim, so it applies to whatever is available to those members. A candidate is so available, and therefore falls under the checks; a private result is not available and does not.
\end{proof}

\begin{proposition}[Review neither necessary nor sufficient]\label{p:review}
Passing peer review is neither necessary nor sufficient for a result to be scientific.
\end{proposition}

\begin{proof}
Not sufficient: review within an enterprise whose central claims are exempt from removal does not make the enterprise self-cleansing, and by Corollary~4.6 of SCI-01-a such an enterprise is not a science, however its journals are organized. Not necessary: by Theorem~4.3 of SCI-01-a, whether a record is self-cleansing depends on whether the refutation schedule for each false claim has no last element, that is, on the checks that follow admission, not on any check that precedes it. A claim admitted without prior review to the record of an enterprise with standing checks is as removable as one admitted after review.
\end{proof}

The proof of Proposition~\ref{p:review} does not make review idle. Review is a check at the point of admission, and it reduces the number of errors admitted. It bears on how much of the record is false at a given time and how quickly the record improves, which is the subject of SCI-19, not on whether the enterprise is a science.

\begin{corollary}[Trust]\label{c:trust}
A candidate does not yet merit the provisional trust that Principle~P3 of SCI-01-a accords to the record. Treating it as established before it has been checked is admitting it to the record by a lower standard, which is permissible only if the checks that follow admission are standing.
\end{corollary}

The proof of Poincaré's conjecture is an instance. Perelman posted his papers to arXiv in 2002 and 2003 and did not submit them to a journal. Several teams then checked and expanded the arguments in detail \citep{KleinerLott2008}, and the result entered the record of mathematics through those checks rather than through review in the journal sense. On the present account the case is not anomalous: the result was a candidate exposed to the standing checks of the discipline, and it was admitted when those checks had been carried out.

\section{Objections and replies}\label{s:objections}

\begin{objection}[Gatekeeping]
If unreviewed results can be scientific, anyone who posts a document can claim the authority of science for it.
\end{objection}

\begin{reply}
Being a candidate confers no authority: by Corollary~\ref{c:trust} authority belongs to the record. Nor is posting sufficient for candidacy. A candidate must be made available to the members of an enterprise in a form that allows them to check it, which requires that the evidence and reasoning be stated, as \citet{Longino1990} argues objectivity requires, and that the enterprise have standing checks that will reach it.
\end{reply}

\begin{objection}[Norms of communication]
Scientific status has traditionally been tied to public communication through recognized channels, as Merton's norm of communism requires \citep{Merton1973}. A preprint server is not such a channel.
\end{objection}

\begin{reply}
What the norm requires is that findings be common property, open to use and criticism. A public preprint satisfies the requirement more fully than a paper behind a journal's paywall. The criterion of SCI-01 explains why the norm matters: communication is the precondition of checks by others.
\end{reply}

\section{Conclusion}\label{s:conclusion}

An unreviewed result is scientific if it is a candidate of a scientific enterprise, exposed to its standing checks; a private result is apt for science but not yet scientific; and no result belongs to the record, and merits its trust, until the enterprise has checked it. Peer review is an important check at admission, but the mark of science lies in the checks that never stop.

\begin{thebibliography}{9}
\bibitem[Ginsparg(2011)]{Ginsparg2011} Ginsparg, P. (2011). ArXiv at 20. \emph{Nature}, 476, 145--147.
\bibitem[Kleiner and Lott(2008)]{KleinerLott2008} Kleiner, B., and Lott, J. (2008). Notes on Perelman's papers. \emph{Geometry and Topology}, 12, 2587--2855.
\bibitem[Longino(1990)]{Longino1990} Longino, H.~E. (1990). \emph{Science as Social Knowledge: Values and Objectivity in Scientific Inquiry}. Princeton University Press.
\bibitem[Merton(1973)]{Merton1973} Merton, R.~K. (1973). The normative structure of science. In \emph{The Sociology of Science: Theoretical and Empirical Investigations}, 267--278. University of Chicago Press. Originally published 1942.
\end{thebibliography}

\end{document}
